% Independent support for the main Part A finite-image/fusion argument.
% Parent assigns statement numbers; all numerical inputs are checked in
% work/ghost_code/check_main_orthogonal_generation.py and its actual output.
\subsection{Two complete orthogonal radicals and the determinant correction}

\begin{lemma}[Generation by the two displayed opposite radicals]
\label{lem:two-orthogonal-radicals}
Let \(V\) be the nondegenerate minus-type quadratic
space of dimension eighteen over \(\mathbb F_3\),
with polar form \(B\) and
\(q(v)=\tfrac12 B(v,v)\).
For a nonzero singular vector \(e\), put
\[
 T_e(v)x=x+B(x,e)v-B(x,v)e-q(v)B(x,e)e,
 \qquad v\in e^\perp,
 \tag{OG.1}
\]
and \(R_e=\{T_e(v):v\in e^\perp\}\).
For any singular \(f\) with \(B(e,f)\ne0\),
\[
                 \langle R_e,R_f\rangle=\Omega^-(18,3).
 \tag{OG.2}
\]
Consequently the actual subgroup generated by the
\(360\) quaternion multiplicity matrices in the
retained certificate is \(\mathrm{SO}^-(18,3)\).
\end{lemma}
\begin{proof}
Direct substitution in (OG.1) shows that it
preserves \(B\), has determinant one, and obeys
\[
 \begin{gathered}
 T_e(v+w)=T_e(v)T_e(w),\qquad
 T_e(v+ae)=T_e(v),\\
 hT_e(v)h^{-1}=T_{he}(hv)\quad(h\in O(V)).
 \end{gathered}
 \tag{OG.3}
\]
Its parameter kernel is precisely \(\langle e\rangle\):
normalize \(B(e,f)=1\), and compare its action
on \(f\) and on \(\langle e,f\rangle^\perp\).
Thus \(R_e\) is the full additive group of
\(e^\perp/\langle e\rangle\), of order \(3^{16}\).
The same formula shows that scaling \(e\)
does not change \(R_e\).

Write \(W=\langle e,f\rangle^\perp\).
Every singular line not orthogonal to \(e\)
has a unique representative
\[
                 f+w-q(w)e\qquad(w\in W),
 \tag{OG.4}
\]
because its pairing with \(e\) can be normalized
to one. Equation (OG.1) takes \(f\) to this
representative. Hence \(R_e\) is regular on
these lines.
Put \(J=\langle R_e,R_f\rangle\).
The space \(W\) is nondegenerate of dimension
sixteen, so it contains a vector \(w\) with
\(q(w)\ne0\): otherwise polarization would
make its form identically zero.
Now \(T_f(w)e=e+w-q(w)f\) is not orthogonal
to \(e\). By (OG.4) an element of \(R_e\)
takes its line to \(\langle f\rangle\).
Therefore the \(J\)-orbit of \(\langle e\rangle\)
contains \(\langle f\rangle\), and then all
singular lines not orthogonal to \(e\).

Let \(u\) be singular, independent of \(e\),
with \(B(u,e)=0\). Nondegeneracy makes the
two linear functionals \(B(e,-),B(u,-)\)
independent, so choose \(a\) with
\(B(e,a)=B(u,a)=1\). The explicit vector
\[
                   z=a-q(a)e
 \tag{OG.5}
\]
satisfies \(q(z)=0\) and
\(B(e,z)=B(u,z)=1\).
Its line is already in the \(J\)-orbit of
\(\langle e\rangle\). Conjugation in (OG.3)
therefore puts \(R_z\) inside \(J\).
Its regular action takes \(\langle e\rangle\)
to \(\langle u\rangle\).
Thus \(J\) is transitive on all singular
lines, and contains every one of their full
radicals by (OG.3).

The matrices (OG.1) are Taylor's Siegel
transformations \(\rho_{e,-v}\).
His Theorems11.24 and11.46 state that these
transformations belong to and generate
\(\Omega(V)\); the exceptional space there
is four-dimensional over \(\mathbb F_2\),
and does not occur here
\cite[pp.~150,161]{GhostTaylorClassical}.
This proves both containments in (OG.2).

For the actual matrices, the independent
certificate verifies the nondegenerate symmetric
form, its minus discriminant, and determinant
one for all \(360\) generators.
It multiplies the \(16\) specified conjugating
words and obtains the matrices (OG.1) with
parameters spanning \(e^\perp/\langle e\rangle\).
The actual generator with zero-based index11
takes \(e\) to a singular vector pairing
nontrivially with it. Hence the actual image
contains both radicals and therefore \(\Omega(V)\).
Direct multiplication of the retained reflection
factorizations gives spinor norm one for
\(336\) generators and the nonsquare class
for \(24\) generators.
Taylor's Theorem11.51 gives the actual quotient
map
\[
 \theta:\mathrm{SO}(V)\longrightarrow
 \mathbb F_3^\times/(\mathbb F_3^\times)^2,
 \qquad \ker\theta=\Omega(V)
 \tag{OG.6}
\]
and its surjectivity
\cite[p.~164]{GhostTaylorClassical}.
The actual subgroup is contained in \(\mathrm{SO}(V)\),
contains this kernel, and meets the other coset.
It is therefore exactly \(\mathrm{SO}^-(18,3)\),
as asserted.
\end{proof}

\begin{lemma}[An orthogonal centralizer element of determinant minus one]
\label{lem:orthogonal-centralizer-determinant}
For the actual subgroup \(H\simeq C_3^2\)
on the eighteen-dimensional multiplicity space,
the orthogonal centralizer \(C_{O(V)}(H)\)
contains an involution of determinant minus one.
Consequently every conjugacy from this \(H\)
to another subgroup obtained inside \(O(V)\)
is also realized inside the actual geometric
image \(\mathrm{SO}^-(18,3)\), with the same
induced map on \(H\).
\end{lemma}
\begin{proof}
Let \(A=\mathbb F_3[H]\), with involution
\(\overline g=g^{-1}\) on its group basis,
augmentation ideal \(\mathfrak m\), and
\(\nu=\sum_{g\in H}g\).
For generators \(s,t\) of \(H\),
\[
 A\simeq\mathbb F_3[X,Y]/(X^3,Y^3),
 \quad X=s-1,quad Y=t-1,quad \nu=X^2Y^2.
 \tag{OG.7}
\]
Its socle is \(\mathbb F_3\nu\).
The actual norm matrix \(\nu|_V\) has rank
two. Choose \(v_1,v_2\) whose norm images
are independent. The \(A\)-linear map
\(A^2\to V\), \((a,b)\mapsto av_1+bv_2\),
is injective on the socle of \(A^2\).
Every nonzero submodule of \(A^2\) meets
that socle: apply a largest nonzero power
of the nilpotent ideal \(\mathfrak m\)
to it. Thus this map has zero kernel.
Both spaces have dimension eighteen over
\(\mathbb F_3\), so it is an isomorphism.
The certificate constructs such a basis
from standard columns1 and2, using zero-based
indices.

The actual form \(B\) is \(H\)-invariant.
Define the \(A\)-valued form
\[
 h(x,y)=\sum_{g\in H}B(x,gy)g.
 \tag{OG.8}
\]
Changing the summation index proves
\(h(ax,y)=a h(x,y)\),
\(h(x,ay)=\bar a h(x,y)\), and
\(h(y,x)=\overline{h(x,y)}\).
The coefficient of the identity in (OG.8)
is exactly \(B(x,y)\).
In an \(A\)-basis the Gram map of \(h\)
is injective: a vector in its kernel pairs
to zero with every vector under \(h\),
and hence under \(B\). Since \(A^2\) is
finite, the Gram map is bijective.
Its reduction modulo \(\mathfrak m\)
is consequently a nondegenerate symmetric
form on \(\mathbb F_3^2\).
Polarization in odd characteristic supplies
a vector with nonzero norm for that form.
Lift it to \(v\in V\). Then
\(a=h(v,v)\) is a self-adjoint unit of
the local ring \(A\).

The explicit \(A\)-linear map
\[
                R(x)=x-2h(x,v)a^{-1}v
 \tag{OG.9}
\]
preserves \(h\): expansion of the two
arguments gives cancelling terms with
coefficients \(-2,-2,4\).
It therefore preserves \(B\) by taking
the identity coefficient, and it centralizes
\(H\) by \(A\)-linearity.
The unit norm gives the actual direct sum
\[
 V=Av\oplus\{x:h(x,v)=0\}.
 \tag{OG.10}
\]
The first summand is free of rank one over
\(A\), since \(cv=0\) implies \(ca=0\).
It has \(\mathbb F_3\)-dimension nine.
The map \(R\) is minus one on this summand
and the identity on the second. Thus it
is an involution with determinant
\((-1)^9=-1\).

These formulas also construct an actual
matrix: the independent certificate forms
(OG.8), inverts \(a\) in the nine-dimensional
group algebra, and evaluates (OG.9).
It checks \(R^2=1\), \(R^TFR=F\),
\(RH_i=H_iR\) for both displayed generators
of \(H\), determinant minus one, and
two eigenspaces of dimension nine.
For an orthogonal conjugator \(c\), retain
\(c\) when its determinant is one and
replace it by \(cR\) otherwise.
Then \(cR\in\mathrm{SO}(V)\), and
\((cR)h(cR)^{-1}=chc^{-1}\) for every
\(h\in H\). The preceding lemma identifies
this special orthogonal group with the
actual geometric image. This proves the
claimed fusion comparison with its exact map.
\end{proof}

The program \texttt{check\_main\_orthogonal\_generation.py}
is separate from the program that generated the
matrices. Its output and JSON record include
the input hashes, the explicit two-radical
orbit bridge, the norm basis, the hermitian
unit and its inverse, and the full determinant
correction matrix. No enumeration of the
orthogonal group is asserted.
